% ============================================================
%  C: Como Programar -- 6ª Edição | Deitel & Deitel
%  Capítulo 3 -- Desenvolvimento Estruturado de Programas em C
%  FAZENDO A DIFERENÇA (3.47, 3.48, 3.49)
% ============================================================
\documentclass[12pt,a4paper]{article}
\usepackage[utf8]{inputenc}
\usepackage[T1]{fontenc}
\usepackage[brazil]{babel}
\usepackage{geometry}
\geometry{top=2.5cm,bottom=2.5cm,left=3cm,right=3cm}
\usepackage{parskip}\setlength{\parskip}{6pt}\setlength{\parindent}{0pt}
\usepackage{xcolor,tcolorbox,enumitem,listings,fancyhdr,titlesec,hyperref,booktabs,array,amsmath}
\tcbuselibrary{skins,breakable}

\definecolor{deitelblue}{RGB}{0,70,127}
\definecolor{deitelgray}{RGB}{240,242,245}
\definecolor{deitelgreen}{RGB}{0,110,60}
\definecolor{deitellightblue}{RGB}{220,235,250}
\definecolor{deitelred}{RGB}{160,20,20}
\definecolor{deitellorange}{RGB}{200,90,0}
\definecolor{ecogreen}{RGB}{0,120,50}
\definecolor{ecolightgreen}{RGB}{210,240,220}
\definecolor{saude}{RGB}{150,0,60}
\definecolor{saudelightred}{RGB}{245,215,225}

\newtcolorbox{boaprática}[1][]{colback=deitellightblue,colframe=deitelblue,
  fonttitle=\bfseries\small,title={Boa prática de programação},breakable,#1}
\newtcolorbox{respostabox}[1][]{colback=deitelgray,colframe=deitelblue!60,
  fonttitle=\bfseries\small\color{deitelblue},title={Resposta detalhada},breakable,#1}
\newtcolorbox{enunciadoBox}[1][]{colback=white,colframe=deitelblue,
  fonttitle=\bfseries,title={#1},breakable}
\newtcolorbox{saida}[1][]{colback=black!88,colframe=deitelblue,
  fontupper=\ttfamily\small\color{green!70},
  title={\small\color{white}Saída do programa},breakable,#1}
\newtcolorbox{pesquisa}[1][]{colback=ecolightgreen,colframe=ecogreen,
  fonttitle=\bfseries\small,title={Pesquisa externa realizada},breakable,#1}
\newtcolorbox{atencao}[1][]{colback=yellow!10,colframe=deitellorange,
  fonttitle=\bfseries\small,title={Atenção},breakable,#1}

\setlist[enumerate,1]{label=\textbf{\alph*)},leftmargin=2em}
\setlist[itemize]{leftmargin=2em}

\lstset{
  basicstyle=\ttfamily\small,backgroundcolor=\color{deitelgray},
  frame=single,framerule=0.4pt,rulecolor=\color{deitelblue!40},
  breaklines=true,showstringspaces=false,
  keywordstyle=\color{deitelblue}\bfseries,
  commentstyle=\color{deitelgreen}\itshape,
  stringstyle=\color{deitelred},
  language=C,numbers=left,numberstyle=\tiny\color{gray},
  xleftmargin=18pt,xrightmargin=5pt,stepnumber=1,
}

\pagestyle{fancy}\fancyhf{}
\fancyhead[L]{\small\color{deitelblue}\textbf{C: Como Programar -- 6ª ed. | Deitel \& Deitel}}
\fancyhead[R]{\small\color{deitelblue}Cap. 3 -- Fazendo a Diferença}
\fancyfoot[C]{\small\thepage}
\renewcommand{\headrulewidth}{0.4pt}

\titleformat{\section}[block]{\large\bfseries\color{deitelblue}}{\thesection.}{0.5em}{}[\titlerule]
\titleformat{\subsection}[block]{\normalsize\bfseries\color{deitelblue!80}}{}{}{}

\hypersetup{colorlinks=true,linkcolor=deitelblue,urlcolor=ecogreen}

\begin{document}

\begin{titlepage}
  \centering\vspace*{2cm}
  {\color{deitelblue}\rule{\linewidth}{2pt}}\\[0.4cm]
  {\huge\bfseries\color{deitelblue}C: Como Programar}\\[0.2cm]
  {\large\color{deitelblue}6ª Edição $\cdot$ Paul Deitel \& Harvey Deitel}\\[0.2cm]
  {\color{deitelblue}\rule{\linewidth}{2pt}}\\[1cm]
  {\LARGE\bfseries Capítulo 3}\\[0.3cm]
  {\Large\color{deitelblue!80}Desenvolvimento Estruturado de Programas em C}\\[0.8cm]
  \begin{tcolorbox}[colback=ecolightgreen,colframe=ecogreen,width=0.85\linewidth]
    \centering{\large\bfseries\color{ecogreen}Fazendo a Diferença}\\[0.3cm]
    {\normalsize Exercícios 3.47 -- 3.49\\[0.2cm]
    Crescimento populacional, frequência cardíaca, criptografia}
  \end{tcolorbox}
  \vfill{\small Abordagem incremental Deitel}
\end{titlepage}

\tableofcontents\newpage

% ============================================================
\section{Exercício 3.47 -- Calculadora de Crescimento Populacional Mundial}
% ============================================================

\begin{enunciadoBox}[Enunciado 3.47]
Use a Web para determinar a população mundial atual e sua taxa de crescimento anual.
Escreva uma aplicação que leia esses valores, depois apresente a população mundial
estimada após um, dois, três, quatro e cinco anos.
\end{enunciadoBox}

\subsection*{Pesquisa: População Mundial e Taxa de Crescimento}

\begin{pesquisa}
Segundo dados das Nações Unidas (UN/DESA), a população mundial em 2024 era de
aproximadamente \textbf{8,1 bilhões de pessoas}, com uma taxa de crescimento anual
de cerca de \textbf{0,87\%} (aproximadamente 70 milhões de pessoas a mais por ano).

\textbf{Fórmula de projeção populacional} (juros compostos aplicada a populações):
\[
  P_n = P_0 \cdot (1 + r)^n
\]
onde $P_n$ é a população após $n$ anos, $P_0$ é a população atual, e $r$ é a taxa
de crescimento anual (em decimal: 0,0087 para 0,87\%).
\end{pesquisa}

\subsection*{Solução em C}

\begin{respostabox}
\begin{lstlisting}
/* Ex3.47: crescimento_populacional.c
   Estima populacao mundial nos proximos 5 anos */
#include <stdio.h>

int main( void )
{
    float populacao;     /* populacao mundial atual */
    float taxa;          /* taxa de crescimento anual (decimal) */
    float projecao;      /* populacao projetada     */
    int anos = 1;        /* contador de anos        */

    printf( "Digite a populacao mundial atual (em bilhoes): " );
    scanf( "%f", &populacao );

    printf( "Digite a taxa anual de crescimento (decimal, "
            "ex: 0.0087 para 0.87%%): " );
    scanf( "%f", &taxa );

    printf( "\n=== PROJECAO POPULACIONAL ===\n" );
    printf( "Ano atual: %.4f bilhoes\n", populacao );

    projecao = populacao;
    while ( anos <= 5 ) {
        projecao = projecao * ( 1.0f + taxa );
        printf( "Apos %d ano(s): %.4f bilhoes\n", anos, projecao );
        ++anos;
    }

    return 0;
} /* fim da funcao main */
\end{lstlisting}

\begin{saida}
Digite a populacao mundial atual (em bilhoes): 8.1\\
Digite a taxa anual de crescimento (decimal, ex: 0.0087 para 0.87\%): 0.0087\\
\\
=== PROJECAO POPULACIONAL ===\\
Ano atual: 8.1000 bilhoes\\
Apos 1 ano(s): 8.1705 bilhoes\\
Apos 2 ano(s): 8.2416 bilhoes\\
Apos 3 ano(s): 8.3133 bilhoes\\
Apos 4 ano(s): 8.3856 bilhoes\\
Apos 5 ano(s): 8.4585 bilhoes
\end{saida}

\medskip\textbf{Reflexão:} O crescimento composto faz com que pequenas taxas anuais
gerem grandes mudanças ao longo do tempo. Em apenas 5 anos, a população mundial
projetada cresce cerca de \textbf{358 milhões de pessoas} -- o equivalente à
população dos EUA! Esse exemplo ilustra a importância de discussões sobre
sustentabilidade, recursos naturais e planejamento global.

\textit{Nota:} O \texttt{\%\%} dentro da string de \texttt{printf} é a forma de
exibir o caractere \texttt{\%} literalmente.
\end{respostabox}

\newpage

% ============================================================
\section{Exercício 3.48 -- Calculadora de Frequência Cardíaca}
% ============================================================

\begin{enunciadoBox}[Enunciado 3.48]
Segundo a American Heart Association (AHA), a fórmula para calcular a frequência
cardíaca máxima é \textbf{220 menos a idade em anos}. A frequência cardíaca-alvo
durante exercícios deve estar entre \textbf{50\% e 85\%} da máxima.
Crie um programa que leia a data de nascimento e a data atual (dia, mês, ano de cada),
calcule e exiba a idade da pessoa, sua frequência cardíaca máxima e a faixa-alvo.
\end{enunciadoBox}

\subsection*{Pesquisa: Cálculo da Idade}

\begin{pesquisa}
\textbf{Cálculo correto da idade em anos completos:}
\begin{enumerate}[noitemsep]
  \item Diferença básica: \texttt{idade = ano\_atual - ano\_nascimento}
  \item Ajuste: se o aniversário ainda \textit{não} ocorreu este ano, subtrair 1
  da idade. O aniversário ocorreu se: \texttt{mês\_atual > mês\_nascimento}, OU
  (\texttt{mês\_atual == mês\_nascimento} E \texttt{dia\_atual >= dia\_nascimento}).
\end{enumerate}

\textbf{Fórmulas da AHA:}
\begin{itemize}[noitemsep]
  \item Frequência máxima: \(\text{FC}_\text{máx} = 220 - \text{idade}\)
  \item Faixa-alvo: 50\% a 85\% da máxima
\end{itemize}
\end{pesquisa}

\subsection*{Solução em C}

\begin{respostabox}
\begin{lstlisting}
/* Ex3.48: frequencia_cardiaca.c
   Calcula idade, FC maxima e faixa-alvo durante exercicios */
#include <stdio.h>

int main( void )
{
    int diaNasc, mesNasc, anoNasc; /* data de nascimento */
    int diaAtual, mesAtual, anoAtual; /* data atual      */
    int idade;
    int fcMaxima;
    float fcMinima, fcAlvo; /* limites inferior e superior em bpm */

    /* Entrada: data de nascimento */
    printf( "=== Calculadora de Frequencia Cardiaca ===\n\n" );
    printf( "DATA DE NASCIMENTO\n" );
    printf( "Dia: " );      scanf( "%d", &diaNasc );
    printf( "Mes: " );      scanf( "%d", &mesNasc );
    printf( "Ano: " );      scanf( "%d", &anoNasc );

    /* Entrada: data atual */
    printf( "\nDATA ATUAL\n" );
    printf( "Dia: " );      scanf( "%d", &diaAtual );
    printf( "Mes: " );      scanf( "%d", &mesAtual );
    printf( "Ano: " );      scanf( "%d", &anoAtual );

    /* Calcula a idade em anos completos */
    idade = anoAtual - anoNasc;
    /* Se o aniversario ainda nao chegou neste ano, subtrai 1 */
    if ( mesAtual < mesNasc ) {
        idade = idade - 1;
    }
    else if ( mesAtual == mesNasc && diaAtual < diaNasc ) {
        idade = idade - 1;
    }

    /* Calcula FC maxima e faixa-alvo */
    fcMaxima = 220 - idade;
    fcMinima = fcMaxima * 0.50f;  /* 50% da maxima */
    fcAlvo   = fcMaxima * 0.85f;  /* 85% da maxima */

    /* Exibe os resultados */
    printf( "\n=== RESULTADO ===\n" );
    printf( "Idade: %d anos\n", idade );
    printf( "Frequencia cardiaca maxima: %d bpm\n", fcMaxima );
    printf( "Faixa-alvo durante exercicios: "
            "%.1f a %.1f bpm\n", fcMinima, fcAlvo );

    printf( "\nATENCAO: Sempre consulte um medico antes de iniciar\n" );
    printf( "         ou modificar um programa de exercicios.\n" );

    return 0;
} /* fim da funcao main */
\end{lstlisting}

\begin{saida}
=== Calculadora de Frequencia Cardiaca ===\\[2pt]
DATA DE NASCIMENTO\\
Dia: 15\\Mes: 8\\Ano: 1995\\[2pt]
DATA ATUAL\\
Dia: 5\\Mes: 5\\Ano: 2026\\[2pt]
=== RESULTADO ===\\
Idade: 30 anos\\
Frequencia cardiaca maxima: 190 bpm\\
Faixa-alvo durante exercicios: 95.0 a 161.5 bpm\\[2pt]
ATENCAO: Sempre consulte um medico antes de iniciar\\
\hphantom{ATENCAO:} ou modificar um programa de exercicios.
\end{saida}
\end{respostabox}

\begin{atencao}
A fórmula 220 -- idade é uma \textit{estimativa simplificada} da AHA. Pesquisas
mais recentes (Tanaka et al., 2001) sugerem fórmulas mais precisas, como
$\text{FC}_\text{máx} = 208 - 0{,}7 \times \text{idade}$. As frequências reais
variam com saúde, condição física e sexo. Consulte sempre um médico.
\end{atencao}

\newpage

% ============================================================
\section{Exercício 3.49 -- Criptografia de 4 Dígitos}
% ============================================================

\begin{enunciadoBox}[Enunciado 3.49]
Uma empresa transmite dados como inteiros de quatro dígitos. Para codificar:
\begin{enumerate}[noitemsep]
  \item Substitua cada dígito por: \((dígito + 7) \% 10\)
  \item Troque o 1º com o 3º dígito; troque o 2º com o 4º dígito.
\end{enumerate}
Escreva um programa que codifique um inteiro de 4 dígitos. Escreva outro programa
que decodifique (revertendo o esquema).
\end{enunciadoBox}

\subsection*{Análise do Algoritmo}

\textbf{Codificação:} se o dígito é $d$, o novo dígito é $(d + 7) \mod 10$.

Exemplo com o número \textbf{1234}:
\begin{enumerate}[noitemsep]
  \item Dígitos originais: $d_1=1, d_2=2, d_3=3, d_4=4$
  \item Após $(d+7) \mod 10$: $d_1=8, d_2=9, d_3=0, d_4=1$
  \item Após troca (1º↔3º, 2º↔4º): $d_1=0, d_2=1, d_3=8, d_4=9$
  \item Número codificado: \textbf{0189}
\end{enumerate}

\textbf{Decodificação:} reverter ambas as operações.
\begin{itemize}[noitemsep]
  \item Inversa da troca: trocar novamente 1º↔3º e 2º↔4º (a troca é sua própria inversa).
  \item Inversa de $(d + 7) \mod 10$: $(d + 3) \mod 10$ (pois $7+3=10\equiv 0$).
\end{itemize}

\subsection*{Programa 1 -- Codificador}

\begin{respostabox}
\begin{lstlisting}
/* Ex3.49(a): codificador.c
   Codifica um inteiro de 4 digitos */
#include <stdio.h>

int main( void )
{
    int numero;
    int d1, d2, d3, d4;       /* digitos originais          */
    int e1, e2, e3, e4;       /* digitos apos (+7) mod 10   */
    int codificado;

    printf( "Digite um inteiro de 4 digitos: " );
    scanf( "%d", &numero );

    /* Extrai os 4 digitos */
    d1 = numero / 1000;
    d2 = ( numero / 100 ) % 10;
    d3 = ( numero / 10 ) % 10;
    d4 = numero % 10;

    /* Aplica (digito + 7) mod 10 */
    e1 = ( d1 + 7 ) % 10;
    e2 = ( d2 + 7 ) % 10;
    e3 = ( d3 + 7 ) % 10;
    e4 = ( d4 + 7 ) % 10;

    /* Troca: 1o <-> 3o, 2o <-> 4o */
    /* O codificado fica na ordem: e3, e4, e1, e2 */
    codificado = e3 * 1000 + e4 * 100 + e1 * 10 + e2;

    printf( "Numero codificado: %04d\n", codificado );

    return 0;
} /* fim da funcao main */
\end{lstlisting}

\textit{Nota:} O especificador \texttt{\%04d} garante que o número seja exibido com
4 dígitos, preenchendo com zeros à esquerda se necessário (ex.: \texttt{189} \(\to\)
\texttt{0189}).
\end{respostabox}

\subsection*{Programa 2 -- Decodificador}

\begin{respostabox}
\begin{lstlisting}
/* Ex3.49(b): decodificador.c
   Decodifica um inteiro de 4 digitos previamente codificado */
#include <stdio.h>

int main( void )
{
    int numero;
    int c1, c2, c3, c4;       /* digitos do numero codificado */
    int d1, d2, d3, d4;       /* digitos originais            */
    int decodificado;

    printf( "Digite um inteiro codificado de 4 digitos: " );
    scanf( "%d", &numero );

    /* Extrai os 4 digitos do codificado */
    c1 = numero / 1000;
    c2 = ( numero / 100 ) % 10;
    c3 = ( numero / 10 ) % 10;
    c4 = numero % 10;

    /* Reverte a troca: 1o <-> 3o, 2o <-> 4o */
    /* Apos reverter: posicoes ficam c3, c4, c1, c2 */

    /* Reverte (d+7) mod 10 aplicando (c+3) mod 10 */
    d1 = ( c3 + 3 ) % 10;
    d2 = ( c4 + 3 ) % 10;
    d3 = ( c1 + 3 ) % 10;
    d4 = ( c2 + 3 ) % 10;

    decodificado = d1 * 1000 + d2 * 100 + d3 * 10 + d4;

    printf( "Numero decodificado: %04d\n", decodificado );

    return 0;
} /* fim da funcao main */
\end{lstlisting}

\begin{saida}
Digite um inteiro codificado de 4 digitos: 189\\
Numero decodificado: 1234
\end{saida}
\end{respostabox}

\subsection*{Verificação Matemática}

\begin{respostabox}
\textbf{Por que $(c + 3) \mod 10$ inverte $(d + 7) \mod 10$?}

Da operação direta: $c = (d + 7) \mod 10$.

Aplicando a inversa: $(c + 3) \mod 10 = ((d + 7) + 3) \mod 10 = (d + 10) \mod 10 = d \mod 10 = d$.

Como $d$ é um dígito (0 a 9), $d \mod 10 = d$. Logo, recuperamos o dígito original.

\medskip\textbf{Verificação completa com 1234 \(\to\) 0189 \(\to\) 1234:}

\textbf{Codificação 1234:}
\begin{itemize}[noitemsep]
  \item Dígitos: 1, 2, 3, 4
  \item Após +7 mod 10: 8, 9, 0, 1
  \item Após troca: 0, 1, 8, 9 \(\to\) número \textbf{0189}
\end{itemize}

\textbf{Decodificação 0189:}
\begin{itemize}[noitemsep]
  \item Dígitos: 0, 1, 8, 9
  \item Reverte troca: 8, 9, 0, 1
  \item Após +3 mod 10: $(8+3)\%10=1$, $(9+3)\%10=2$, $(0+3)\%10=3$, $(1+3)\%10=4$
  \item Resultado: 1, 2, 3, 4 \(\to\) número \textbf{1234} \checkmark
\end{itemize}
\end{respostabox}

\subsection*{Reflexão sobre Criptografia}

\begin{tcolorbox}[colback=ecolightgreen,colframe=ecogreen,
  title={\bfseries Da brincadeira didática à criptografia real}]

O esquema deste exercício é uma forma muito simples de criptografia, conhecida como
\textbf{cifra de substituição com transposição}. Ele tem grandes limitações:

\begin{itemize}[noitemsep]
  \item \textbf{Reversibilidade óbvia:} Quem conhece o algoritmo decodifica trivialmente.
  \item \textbf{Sem chave secreta:} Não há ``segredo'' que separe usuários legítimos
  de atacantes.
  \item \textbf{Espaço minúsculo de mensagens:} Apenas 10.000 valores possíveis (0--9999).
\end{itemize}

\textbf{Sistemas modernos} usam:
\begin{itemize}[noitemsep]
  \item \textbf{Criptografia simétrica} (AES) -- chave única compartilhada
  \item \textbf{Criptografia de chave pública} (RSA, PGP) -- chave pública para
  cifrar e chave privada para decifrar; baseada em problemas matemáticos
  computacionalmente difíceis (fatoração de inteiros grandes, logaritmo discreto)
  \item \textbf{Funções hash criptográficas} (SHA-256) -- garantem integridade dos dados
\end{itemize}

Esse exercício é uma porta de entrada conceitual para um campo fundamental da
ciência da computação contemporânea: a \textbf{segurança da informação}.
\end{tcolorbox}

\newpage

% ============================================================
\section*{Resumo -- Fazendo a Diferença no Capítulo~3}

\begin{tcolorbox}[colback=deitellightblue,colframe=deitelblue,
  title={\bfseries Os três exercícios e suas conexões}]

\textbf{3.47 -- Crescimento Populacional:}
Aplicação de cálculos compostos em escala global. O laço \texttt{while} permite projetar
cenários ano a ano. Conecta-se ao Ex.~1.10 (carbono) e ao Ex.~2.33 (transporte solidário) --
todos lidam com sustentabilidade e impacto humano no planeta.

\medskip\textbf{3.48 -- Frequência Cardíaca:}
Aplicação direta de \texttt{if...else} para o cálculo cuidadoso da idade,
e da estrutura sequencial para determinar limites cardiovasculares seguros. Conecta-se
ao Ex.~1.11 e ao Ex.~2.32 (IMC) -- a programação como ferramenta para a saúde.

\medskip\textbf{3.49 -- Criptografia:}
Manipulação avançada de dígitos com aritmética modular. Conecta-se aos Exercícios
2.30 (separar dígitos) e 3.35 (palíndromo) -- a aritmética inteira como ferramenta
para resolver problemas surpreendentemente complexos.

\end{tcolorbox}

\begin{boaprática}
  Os três exercícios deste capítulo demonstram como, com apenas as ferramentas dos
  Capítulos~1, 2 e 3 -- variáveis, \texttt{if/else}, \texttt{while} e operações
  aritméticas -- já é possível abordar problemas reais e relevantes. À medida que
  novos recursos da linguagem forem introduzidos nos próximos capítulos, sua
  capacidade de resolver problemas crescerá exponencialmente.
\end{boaprática}

\end{document}
